\section*{Chapitre \ref{ch:g6:life-through-seasons} --- La vie au fil des saisons}

\begin{solution}{exo:g6:life-through-seasons:1}
La durée du jour (la \omterm{def:g6:exploring-our-environment:conditions}{lumière}), la \omterm{def:g6:exploring-our-environment:conditions}{température}, et l'\omterm{def:g6:exploring-our-environment:conditions}{humidité} du sol
ou de l'air, qui suit la pluie et la sécheresse.
\end{solution}

\begin{solution}{exo:g6:life-through-seasons:2}
\omterm{def:g6:life-through-seasons:trees}{Caduc} --- perd toutes ses \omterm{def:g1:plants-around-us:parts}{feuilles} en automne~: chêne, hêtre,
marronnier. \omterm{def:g6:life-through-seasons:trees}{Persistant} --- garde tout l'hiver des \omterm{def:g1:plants-around-us:parts}{feuilles} qui
travaillent~: pin, houx.
\end{solution}

\begin{solution}{exo:g6:life-through-seasons:3}
Une minuscule pousse complète --- des \omterm{def:g1:plants-around-us:parts}{feuilles} miniatures pliées,
enveloppées d'\omterm{def:g1:animals-around-us:covering}{écailles} à l'épreuve du temps. Il a été bâti l'été
précédent.
\end{solution}

\begin{solution}{exo:g6:life-through-seasons:4}
L'annuelle meurt entièrement et passe l'hiver sous forme de \omterm{prop:g3:flowers-fruits-seeds:transform}{graines}.
La partie aérienne de la jonquille meurt, tandis que la \omterm{def:g1:plants-around-us:plant}{plante}
attend sous terre sous forme de \omterm{ex:g4:plants-without-seeds:bulbtuber}{bulbe}.
\end{solution}

\begin{solution}{exo:g6:life-through-seasons:5}
Parce que les pousses de ses \omterm{def:g6:life-through-seasons:buds}{bourgeons} étaient déjà finies l'été
dernier~; il ne manquait au rameau que la chaleur, que l'eau à
l'intérieur lui fournit des semaines avant le jardin.
\end{solution}

\begin{solution}{exo:g6:life-through-seasons:6}
\omterm{def:g3:animals-through-seasons:migration}{Migrer}~: l'hirondelle. \omterm{def:g3:animals-through-seasons:hibernation}{Hiberner}~: le hérisson (ou l'escargot
scellé). Rester actif~: le renard, la pie, le rouge-gorge. Hiverner
sous des formes jeunes --- \omterm{def:g2:animals-grow-up:born}{œuf}, larve ou \omterm{ex:g2:animals-grow-up:butterfly}{chrysalide}~: la plupart des
\omterm{prop:g3:sorting-living-things:groups}{insectes}, comme la sauterelle ou le papillon.
\end{solution}

\begin{solution}{exo:g6:life-through-seasons:7}
Papillon~: une \omterm{ex:g2:animals-grow-up:butterfly}{chrysalide} sous une pierre du chaperon. Sauterelle~:
des \omterm{def:g2:animals-grow-up:born}{œufs} dans le sol. Coccinelle~: un adulte endormi, serré au fond
de la litière. Escargot~: scellé dans sa \omterm{def:g1:animals-around-us:covering}{coquille} dans le talus.
Hirondelle~: en Afrique, sous sa forme migrée.
\end{solution}

\begin{solution}{exo:g6:life-through-seasons:8}
Pour que toute différence entre les quatre tableaux vienne des
saisons et non du relevé --- la règle de l'expérience honnête de
\cref{met:g6:exploring-our-environment:survey}, appliquée au temps
plutôt qu'aux sites.
\end{solution}

\begin{solution}{exo:g6:life-through-seasons:9}
D'abord, beaucoup d'«~absents~» sont présents sous des formes
cachées --- \omterm{def:g2:animals-grow-up:born}{œufs}, \omterm{ex:g2:animals-grow-up:butterfly}{chrysalides}, adultes endormis, \omterm{ex:g4:plants-without-seeds:bulbtuber}{bulbes} --- qu'un
relevé fait en se promenant manque sans creuser ni soulever.
Ensuite, certains ne sont pas perdus mais ailleurs~: les \omterm{def:g3:animals-through-seasons:migration}{migrateurs}
reviendront. La liste de janvier sous-estime les résidents et compte
les voyageurs comme disparus.
\end{solution}

\begin{solution}{exo:g6:life-through-seasons:10}
La \omterm{def:g1:plants-around-us:plant}{plante} est dans son état de réserve souterraine --- vivante, en
attente sous forme de \omterm{ex:g4:plants-without-seeds:bulbtuber}{bulbes}. En mars, les \omterm{ex:g4:plants-without-seeds:bulbtuber}{bulbes} redonneront
\omterm{def:g1:plants-around-us:parts}{feuilles} et \omterm{def:g3:flowers-fruits-seeds:parts}{fleurs}, plus denses si les touffes se sont multipliées.
\end{solution}

\begin{solution}{exo:g6:life-through-seasons:11}
Disparues où, sous quelle forme~? Les grenouilles hivernent
engourdies --- enfouies dans la vase du fond de la mare ou dans des
abris humides à terre --- et passent les mois froids sans activité
jusqu'à ce que le printemps les rappelle à l'eau pour se reproduire.
\end{solution}

\begin{solution}{exo:g6:life-through-seasons:12}
Le problème de l'hiver, pour une \omterm{def:g1:plants-around-us:parts}{feuille} qui travaille, est de
retenir l'eau~: le sol gelé n'en donne presque pas aux \omterm{def:g1:plants-around-us:parts}{racines},
tandis que le vent et le soleil continuent d'en tirer --- le
problème de dessèchement du cloporte, tout l'hiver durant. La cire
et la forme en aiguille de la \omterm{def:g1:plants-around-us:parts}{feuille} persistante ramènent ces
pertes à ce que l'arbre peut se permettre. Une \omterm{def:g1:plants-around-us:parts}{feuille} large et
mince ne peut pas être aussi bien étanchée~; la réponse meilleur
marché de l'arbre \omterm{def:g6:life-through-seasons:trees}{caduc} est de passer ses \omterm{def:g1:plants-around-us:parts}{feuilles} minces par pertes
et profits chaque automne, et de ranger le printemps dans des
\omterm{def:g6:life-through-seasons:buds}{bourgeons}.
\end{solution}

\begin{solution}{pb:g6:life-through-seasons:1}
\textbf{1.} Les \omterm{def:g6:life-through-seasons:buds}{bourgeons} de chaque rameau --- le rangement d'hiver
qui contient les pousses du printemps suivant, bâties l'été d'avant.

\textbf{2.} Entre les images de mars et d'avril (avec les jours qui
s'allongent et les matins qui s'adoucissent dans ses notes).
L'arbre répond à la hausse de \omterm{def:g6:exploring-our-environment:conditions}{température} et à l'allongement de la
\omterm{def:g6:exploring-our-environment:conditions}{lumière}.

\textbf{3.} Des \omterm{prop:g3:sorting-living-things:groups}{insectes} pollinisateurs --- des abeilles avant tout
--- doivent visiter les \omterm{def:g3:flowers-fruits-seeds:parts}{fleurs}, pour que du \omterm{def:g3:flowers-fruits-seeds:parts}{pollen} atteigne les
\omterm{def:g3:flowers-fruits-seeds:parts}{pistils} et que la \omterm{def:g5:flower-to-fruit:steps}{fécondation} suive. Un marron est la \omterm{def:g2:from-seed-to-plant:seed}{graine} formée
à partir d'un ovule fécondé, dans sa bogue épineuse.

\textbf{4.} \omterm{def:g6:life-through-seasons:trees}{Caduc}. Son printemps suivant traverse les images
d'hiver sous forme de \omterm{def:g6:life-through-seasons:buds}{bourgeons} sur les rameaux nus.

\textbf{5.} Sous la pelouse, sous forme de \omterm{ex:g4:plants-without-seeds:bulbtuber}{bulbes} --- la réserve
souterraine de la vivace.

\textbf{6.} Sous forme de \omterm{prop:g3:flowers-fruits-seeds:transform}{graines} dans le sol~: la stratégie de
l'annuelle --- les \omterm{def:g1:plants-around-us:plant}{plantes} elles-mêmes sont mortes en automne.

\textbf{7.} L'hirondelle~: en Afrique, sous sa forme migrée, à la
suite des \omterm{prop:g3:sorting-living-things:groups}{insectes} volants. Les abeilles~: dans la ruche, en grappe
d'hiver, au chaud autour de leur reine.

\textbf{8.} Rester actif. La pie doit trouver à manger pendant les
mois les plus maigres --- le problème même que l'hirondelle évite en
partant.

\textbf{9.} Par exemple~: les \omterm{def:g6:life-through-seasons:buds}{bourgeons} éclatent (mars--avril) quand
les jours passent de $11$ à $13$ heures et que les matins dépassent
\qty{8}{\degreeCelsius}~; les \omterm{def:g1:plants-around-us:parts}{feuilles} dorent et tombent
(octobre--novembre) quand les jours raccourcissent et que les matins
fraîchissent. Les événements de l'arbre s'associent à des
changements de \omterm{def:g6:exploring-our-environment:conditions}{conditions}, image après image.

\textbf{10.} Le carré n'est pas mort~: il s'est retiré dans sa
réserve --- vivant, sous forme de \omterm{ex:g4:plants-without-seeds:bulbtuber}{bulbes}, au même endroit.
«~Disparu~» mérite toujours~: disparu \emph{où}, sous quelle
\emph{forme}~?

\textbf{11.} Des \omterm{def:g1:plants-around-us:tree}{branches} nues portant des \omterm{def:g6:life-through-seasons:buds}{bourgeons}~; une pie
quelque part, probablement~; en dessous, une litière de \omterm{def:g1:plants-around-us:parts}{feuilles} et
le sol nu des jonquilles qui cache les \omterm{ex:g4:plants-without-seeds:bulbtuber}{bulbes}~; des \omterm{prop:g3:flowers-fruits-seeds:transform}{graines} de
coquelicot dans la terre du chemin. Les paris sûrs sont les
\omterm{def:g6:life-through-seasons:buds}{bourgeons} et la nudité --- ils découlent de la machinerie caduque de
l'arbre et d'une année de preuves~; la pie est probable, non
promise.

\textbf{12.} Cela montre que la chaleur était la condition
manquante --- la \omterm{def:g6:exploring-our-environment:conditions}{lumière} à l'intérieur en février n'est pas plus
longue que dehors, mais l'eau et la pièce sont tièdes, et cela a
suffi. Expérience honnête~: deux rameaux du même arbre, tous deux
dans l'eau à la même fenêtre, l'un dans une pièce chaude, l'autre
dans un couloir froid --- une seule différence, la chaleur~; si seul
le rameau au chaud débourre, le verdict est acquis.
\end{solution}
